\section{Diagnostics for the old attempted reductions}
\label{sec:p147-diagnostics}

These calculations are independent of the candidate maximality proof.

\subsection{Singularity data do not specify a normalized representative}

The function \(V_\theta-1\) has the same singularity type and the same
multiplier ideals at every positive real weight as \(V_\theta\).
Nevertheless
\[
 P_\theta[V_\theta-1]
   =P_\theta[V_\theta-1]_{\mathcal I}
   =V_\theta\ne V_\theta-1.
\]
Indeed \(V_\theta\) is admissible for both envelopes, while every
normalized \(\theta\)-psh competitor is at most \(V_\theta\).
Their nonpluripolar Monge--Amp\`ere measures are also equal, since
adding a constant does not change the curvature current.
Thus analytic singularities, integral closure, agreement of all
multiplier ideals, or agreement of mass do not by themselves prove that
a particular section-defined representative is maximal.
This is a diagnostic example, not a filtration counterexample.

Similarly, if \(|u-v|\leq C\) globally, their multiplier ideals are
identical at each fixed real weight \(t\): the two integrands differ by
factors between \(\exp(-tC)\) and \(\exp(tC)\).
Consequently an actual bounded \(o(1)\) difference cannot create different
multiplier ideals in the way asserted in the old dossier.
Convergence only away from a base locus, or convergence while changing
the background class, is a different statement.

\subsection{The blowup calculation with its exceptional factor}

Let \(\mu:S=\operatorname{Bl}_p\mathbf P^2\to\mathbf P^2\),
let \(E\) be the exceptional curve, and let
\[
                  L'=\mu^*\mathcal O_{\mathbf P^2}(d)
                                       \otimes\mathcal O_S(E),\qquad d\geq1.
\]
For every positive integer \(m\), the correct section-space identity is
\begin{equation}\label{eq:p147-blowup-sections}
 H^0(S,(L')^{\otimes m})
       =s_E^m\,\mu^*H^0(\mathbf P^2,\mathcal O_{\mathbf P^2}(md)).
\end{equation}
There is no imposed order-\(m\) vanishing condition on the downstairs
polynomial.
To prove the identity, tensor the divisor exact sequence by
\(\mu^*\mathcal O(md)\).  Its restriction term for the \(j\)-th exceptional
factor is \(\mathcal O_E(-j)\), since
\(\mu^*\mathcal O(md)|_E\) is trivial and
\(\mathcal O(E)|_E=\mathcal O_E(-1)\).
It has no sections when \(j\geq1\).  Successively removing \(m\) such
factors shows that all sections are divisible by \(s_E^m\).
Finally \(H^0(S,\mu^*\mathcal O(md))=H^0(\mathbf P^2,\mathcal O(md))\)
because \(\mu_*\mathcal O_S=\mathcal O_{\mathbf P^2}\).
This proves \eqref{eq:p147-blowup-sections} and gives
\[
                     h^0(S,(L')^{\otimes m})=\binom{md+2}{2}.
\]
For a nonzero downstairs polynomial \(t\),
\[
            \operatorname{ord}_E(s_E^m\mu^*t)
                            =m+\operatorname{mult}_p(t).
\]
Thus the \(\operatorname{ord}_E\)-filtration at weight \(m\tau\)
corresponds to downstairs polynomials of degree \(md\) vanishing at \(p\)
to order at least
\[
                    \ell_m=\max\{0,\lceil m(\tau-1)\rceil\}.
\]
For \(0\leq\ell_m\leq md\), choose coordinates
\(p=[1:0:0]\) and count monomials
\(x_0^{md-a-b}x_1^a x_2^b\).
Exactly \(\ell_m(\ell_m+1)/2\) monomials have \(a+b<\ell_m\).
The dimension of the filtered piece is therefore
\[
                 \binom{md+2}{2}-\binom{\ell_m+1}{2}.
\]
It is zero if \(\ell_m>md\).
The normalized maximal weight is \(c=d+1\), and the graded volumes are
\begin{equation}\label{eq:p147-blowup-volumes}
 \operatorname{vol}(W^\tau)=
 \begin{cases}
 d^2,&\tau\leq1,\\
 d^2-(\tau-1)^2,&1<\tau<d+1,\\
 0,&\tau=d+1.
 \end{cases}
\end{equation}
At the terminal weight the dimension is \(md+1\), not quadratic in \(m\).
All strict interior weights have positive volume.

The elementary intersection numbers \(H^2=1\), \(H\cdot E=0\),
\(E^2=-1\), where \(H=\mu^*\mathcal O(1)\), show
\(L'\cdot E=-1\); its Zariski decomposition is \(dH+E\).
This confirms the fixed exceptional factor but, by itself, does not
calculate the normalized envelope of any nontrivial filtered slice.
In particular the old argument did not refute all counterexamples of
this birational type merely by matching an exceptional Lelong number.

Finally, the augmented base locus is not the locus on which every
section vanishes.  For \(H=\mu^*\mathcal O(1)\), the bundle is globally
generated, whereas \(E\subseteq\mathbf B_+(H)\).
One can see the latter directly: for any ample divisor \(B\) on \(S\),
\((H-\varepsilon B)\cdot E<0\), so every effective multiple of
\(H-\varepsilon B\) contains \(E\).
This disproves the old dossier's asserted base-locus obstruction.

